\section{September 28, 2026}
% Based on pages 1--3 of 5510928.pdf.

We revisit stability of Gauss and Radau methods, interpret their
stability functions as Pad\'e approximations to the exponential,
and study reversibility under a change in the direction of time.
We retain \(A\) for the matrix in a linear ODE,
\(\mathsf A=(a_{ij})\) for the RK coefficient matrix, and
\(\tau=t_{n+1}-t_n\). As in the earlier lectures, \(\Phi\) is
the exact flow and \(\Psi\) is the numerical flow.

\subsection{Linear stability and stiff decay}

For the constant-coefficient system \(x'=Ax\), the exact and RK
updates are
\[
 x(t_{n+1})=e^{\tau A}x(t_n),\qquad
 x_{n+1}=R(\tau A)x_n.
\]
Here \(R\) is the rational stability function
\begin{equation}\label{eq:rk-stability-sep-28}
 R(z)=1+z b^{\mathsf T}(I-z\mathsf A)^{-1}\mathbf1,
 \qquad z=\tau\lambda.
\end{equation}
If \(R=P/Q\), its matrix evaluation means
\(R(\tau A)=Q(\tau A)^{-1}P(\tau A)\), whenever defined;
the RK stage system must also be solvable.

Recall the distinction between the strict stability region
\(\mathcal S\), where \(|R(z)|<1\), and the non-strict domain
\(\mathcal D\), where \(|R(z)|\leq1\). In both cases we require
the stages to be defined. The handwritten notes use the non-strict
convention in this lecture. An \(A\)-stable method has
\[
 \{z\in\bb C:\operatorname{Re}z\leq0\}\subseteq\mathcal D.
\]
Gauss and Radau IIA methods are \(A\)-stable for every number
of stages \(s\geq1\).

For the differential equation, asymptotic stability of the origin
means that every solution tends to zero, and is equivalent to
\begin{equation}\label{eq:linear-decay-sep-28}
 \operatorname{Re}\lambda<0\qquad\text{for all }\lambda\in\sigma(A).
\end{equation}
The sign is negative. Mere boundedness allows eigenvalues on the
imaginary axis, provided they are semisimple, as discussed on
September 11. Likewise, at a fixed positive step size,
\[
 R(\tau A)^n\longrightarrow0
 \quad\Longleftrightarrow\quad
 |R(\tau\lambda)|<1\quad\text{for all }\lambda\in\sigma(A).
\]
Gauss and Radau IIA satisfy this strict inequality in the open
left half-plane, so both reproduce asymptotic decay for every
\(\tau>0\) when \eqref{eq:linear-decay-sep-28} holds.

\begin{definition}[\(L\)-stability]
 An \(A\)-stable method is \emph{\(L\)-stable} if
 \[
  R(z)\longrightarrow0
  \quad\text{as }|z|\to\infty\text{ in the left half-plane}.
 \]
\end{definition}

In particular, for a fixed \(\lambda\) with
\(\operatorname{Re}\lambda<0\), an \(L\)-stable method satisfies
\(R(\tau\lambda)\to0\) as \(\tau\to\infty\), matching the
limit of \(e^{\tau\lambda}\). This is a one-step damping property
for very stiff modes, in addition to decay over many fixed steps.
Radau IIA methods are \(L\)-stable; Gauss methods are not.

\subsection{Pad\'e approximation of the exponential}

\begin{definition}[Pad\'e approximant]
 The \([m/n]\) Pad\'e approximant of \(e^z\) is a rational
 function \(P_m(z)/Q_n(z)\), with
 \[
  \deg P_m\leq m,\qquad \deg Q_n\leq n,\qquad Q_n(0)=1,
 \]
 whose Taylor expansion agrees with that of \(e^z\) through
 degree \(m+n\):
 \begin{equation}\label{eq:pade-matching-sep-28}
  Q_n(z)e^z-P_m(z)=O(z^{m+n+1}).
 \end{equation}
\end{definition}

An RK method of order \(p\), applied to the scalar test equation,
necessarily satisfies
\[
 R(z)=e^z+O(z^{p+1}).
\]
This explains the connection between RK stability functions and
rational approximations to the exponential. The Pad\'e property
requires the particular degree and matching conditions above;
it does not hold for every RK stability function.

For an \(s\)-stage RK method, the determinant identity gives
\[
 R(z)=
 \frac{\det(I-z\mathsf A+z\mathbf1 b^{\mathsf T})}
      {\det(I-z\mathsf A)}.
\]
Thus the numerator and denominator have degrees at most \(s\).
A Gauss method has order \(2s\), so its stability function is
the diagonal \([s/s]\) Pad\'e approximant of \(e^z\).
The first two examples are
\begin{align*}
 R_{\mathrm{G},1}(z)&=\frac{1+z/2}{1-z/2},
 &R_{\mathrm{G},1}(z)-e^z&=O(z^3),\\
 R_{\mathrm{G},2}(z)&=\frac{1+z/2+z^2/12}{1-z/2+z^2/12},
 &R_{\mathrm{G},2}(z)-e^z&=O(z^5).
\end{align*}
Consequently, on a linear system these methods give rational
matrix approximations
\[
 e^{\tau A}\approx R(\tau A).
\]
For example, the one-stage Gauss method replaces the exponential by
\((I-\tau A/2)^{-1}(I+\tau A/2)\).

The diagonal exponential approximant can be written as
\begin{equation}\label{eq:gauss-pade-sep-28}
 R_{\mathrm{G},s}(z)=\frac{P_s(z)}{P_s(-z)},
 \qquad
 P_s(z)=\sum_{k=0}^s
 \frac{(2s-k)!\,s!}{(2s)!\,(s-k)!\,k!}\,z^k.
\end{equation}
Its numerator and denominator have the same degree, and their
leading coefficients differ by \((-1)^s\). Hence
\[
 \lim_{|z|\to\infty}R_{\mathrm{G},s}(z)=(-1)^s,
\]
which explains the failure of \(L\)-stability.

\subsubsection{The nonlinear flow operator}

For an autonomous equation \(x'=f(x)\) in \(\bb R^d\), introduce
the differential operator
\[
 \mathcal L g(x)=\sum_{j=1}^d f_j(x)\frac{\partial g}{\partial x_j}(x)
               =Dg(x)f(x).
\]
Along the exact solution,
\[
 \frac{d}{dt}g(x(t))=(\mathcal Lg)(x(t)).
\]
Thus the action of the exact flow on functions is formally
\[
 g(\Phi^\tau(x))=(e^{\tau\mathcal L}g)(x)
 =g(x)+\tau\mathcal Lg(x)
       +\frac{\tau^2}{2}\mathcal L^2g(x)+\cdots.
\]
For smooth functions this can be read as a finite Taylor expansion
with a remainder; convergence of the infinite series needs further
assumptions. This is the operator analogy with the matrix
exponential mentioned in class. For general nonlinear \(f\), the
RK map is not in general obtained by simply substituting
\(\tau\mathcal L\) into its scalar stability function. The formula
\(R(\tau A)x\) is exact for the numerical update of a linear system.

\subsection{Why Radau IIA is \texorpdfstring{\(L\)}{L}-stable}

Suppose \(\mathsf A\) is invertible. From
\eqref{eq:rk-stability-sep-28},
\[
 z(I-z\mathsf A)^{-1}=-\mathsf A^{-1}+O(|z|^{-1}),
\]
and therefore
\begin{equation}\label{eq:rk-limit-sep-28}
 \lim_{|z|\to\infty}R(z)
   =1-b^{\mathsf T}\mathsf A^{-1}\mathbf1.
\end{equation}
If \(b^{\mathsf T}\) coincides with the last row of
\(\mathsf A\), then
\[
 b^{\mathsf T}=e_s^{\mathsf T}\mathsf A
 \quad\Longrightarrow\quad
 1-b^{\mathsf T}\mathsf A^{-1}\mathbf1
 =1-e_s^{\mathsf T}\mathbf1=0.
\]
Here \(e_s\) is the last coordinate vector. Together with
\(A\)-stability, this proves \(L\)-stability.

For collocation, the coefficients are
\[
 b_j=\int_0^1\ell_j(\theta)\,d\theta,\qquad
 a_{ij}=\int_0^{c_i}\ell_j(\theta)\,d\theta.
\]
Right Radau nodes include \(c_s=1\), so \(a_{sj}=b_j\).
This is the stiff accuracy property from September 25, and the
last stage equals \(x_{n+1}\). The nodes are distinct and nonzero,
so \(\mathsf A\) is invertible by the polynomial argument given
there. Since Radau IIA is also \(A\)-stable, every member of the
family is \(L\)-stable.

For the first two members,
\[
 R_{\mathrm{R},1}(z)=\frac1{1-z},\qquad
 R_{\mathrm{R},2}(z)=\frac{1+z/3}{1-2z/3+z^2/6}.
\]
Both tend to zero as \(|z|\to\infty\). More generally, the
denominator has degree \(s\), and the zero limit forces the
numerator to have degree at most \(s-1\). Since Radau IIA has
order \(2s-1\), its stability function is the \([s-1/s]\)
Pad\'e approximant of the exponential.

\subsection{Time reversibility of a one-step method}

For the exact flow of \(x'=f(t,x)\), uniqueness implies
\[
 \Phi^{t_2,t_1}(x(t_1))=x(t_2),\qquad
 \Phi^{t_1,t_2}(x(t_2))=x(t_1).
\]
Whenever both directions of the flow exist,
\begin{equation}\label{eq:exact-reversibility-sep-28}
 \Phi^{t_1,t_2}\circ\Phi^{t_2,t_1}=I.
\end{equation}
Advancing along a solution and then returning along the same
solution recovers the initial value.

\begin{definition}
 A one-step method is \emph{time-reversible}, or
 \emph{symmetric}, if
 \begin{equation}\label{eq:numerical-reversibility-sep-28}
  \Psi^{t_n,t_{n+1}}\circ\Psi^{t_{n+1},t_n}=I.
 \end{equation}
 The backward map uses the same method with step \(-\tau\).
 For an autonomous equation this reads
 \(\Psi^{-\tau}\circ\Psi^\tau=I\).
\end{definition}

For implicit methods, this statement is understood where the
forward and backward steps are well defined, using the locally
unique solution branches of their stage equations.

\begin{example}[Forward Euler is not reversible]
 A forward step gives
 \[
  x_{n+1}=x_n+\tau f(t_n,x_n).
 \]
 Applying forward Euler with step \(-\tau\) from \(t_{n+1}\)
 gives a returned value
 \begin{align*}
  \widehat x_n
   &=x_{n+1}-\tau f(t_{n+1},x_{n+1})\\
   &=x_n+\tau f(t_n,x_n)-\tau f(t_{n+1},x_{n+1}),
 \end{align*}
 which generally differs from \(x_n\). On \(x'=\lambda x\),
 the forward and backward multipliers give
 \[
  \widehat x_n=(1-z)(1+z)x_n=(1-z^2)x_n.
 \]
\end{example}

\subsection{Reversibility of Gauss collocation}

\begin{theorem}
 Gauss collocation methods are time-reversible wherever the
 relevant collocation solutions are uniquely defined.
\end{theorem}

\begin{proof}
 Let \(t_{n+1}=t_n+\tau\), and let \(u\in(\mathcal P_s)^d\)
 be the forward collocation polynomial:
 \[
  u(t_n)=x_n,\qquad
  u'(t_n+c_i\tau)=f(t_n+c_i\tau,u(t_n+c_i\tau)),
 \]
 with \(x_{n+1}=u(t_{n+1})\). The backward step seeks
 \(v\in(\mathcal P_s)^d\) satisfying
 \[
  v(t_{n+1})=x_{n+1},\qquad
  v'(t_{n+1}-c_i\tau)
    =f(t_{n+1}-c_i\tau,v(t_{n+1}-c_i\tau)).
 \]
 Gauss nodes obey \(c_{s+1-i}=1-c_i\). Consequently,
 \[
  t_{n+1}-c_i\tau
    =t_n+(1-c_i)\tau=t_n+c_{s+1-i}\tau.
 \]
 The backward collocation points are precisely the forward
 collocation points in reverse order. The polynomial \(u\)
 therefore satisfies all the backward conditions as well.
 Uniqueness gives \(v=u\), and the returned value is
 \(v(t_n)=u(t_n)=x_n\), proving
 \eqref{eq:numerical-reversibility-sep-28}.
\end{proof}

The derivative in this proof is taken with respect to physical
time; the backward step changes the node locations through
\(-\tau\). Equivalently, in the normalized variable,
\(V(\theta)=U(1-\theta)\) has
\(V'(\theta)=-U'(1-\theta)\), which supplies the negative step
in the backward collocation equations.

On the scalar test equation, reversibility requires
\[
 R(-z)R(z)=1
\]
where both steps are defined. The Gauss Pad\'e formula
\eqref{eq:gauss-pade-sep-28} has exactly this property.
For the one-stage method, it is also immediate from
\[
 x_{n+1}=x_n+\tau f\left(t_n+\frac\tau2,
                              \frac{x_n+x_{n+1}}2\right):
\]
interchanging the endpoints and replacing \(\tau\) by
\(-\tau\) leaves the equation unchanged. The collocation proof
extends this symmetry to every Gauss method and to nonlinear,
time-dependent vector fields.
